\documentclass[11pt]{article}
\usepackage{graphicx}
\usepackage{amssymb}
\usepackage{amsmath}
\usepackage{epstopdf}
\DeclareGraphicsRule{.tif}{png}{.png}{`convert #1 `dirname #1`/`basename #1 .tif`.png}

\textwidth = 6.5 in
\textheight = 9 in
\oddsidemargin = 0.0 in
\evensidemargin = 0.0 in
\topmargin = 0.0 in
\headheight = 0.0 in
\headsep = 0.0 in
\parskip = 0.2in
\parindent = 0.0in

\newtheorem{theorem}{Theorem}
\newtheorem{corollary}[theorem]{Corollary}
\newtheorem{definition}{Definition}

\newcommand{\cut}[1]{}

\title{Matrix Calculus Bootcamp}
\author{Terence Parr and Jeremy Howard}
\begin{document}
\maketitle

\section{Introduction}

three-dimensional coordinates in room yielding a temperature. then can yield temperature and noise or humidity or light.  Sheldon Cooper: you're in my spot.

$\Delta x$ gives what $\Delta y$? rise/run to determine derivative.

neural net equation. optimizing it and calculus comes lurching back into your life like distant relatives around the holidays.

 lots of good cheat sheets like https://www.math.uwaterloo.ca/~hwolkowi/matrixcookbook.pdf and slides like https://www.comp.nus.edu.sg/~cs5240/lecture/matrix-differentiation.pdf
 
given any particular formula, you can reduce it to its scalar form and use the typical scalar derivative to come up with the formula. That can be tedious and error-prone. It's better to learn the matrix/vector derivatives and apply transformations to get the derivative.

\section{Review: Scalar derivative rules}

$
\begin{array}{lcc}
Rule & \text{Scalar function } f(x) & \text{Scalar derivative with repsect to } x\\
\hline
\text{Multiplication by constant} &	cf(x)	&cf'(x)\\
\text{Power Rule}	& x^n	& nx^{n-1}\\
\text{Sum Rule}	& f + g	& f' + g'\\
\text{Difference Rule}	& f - g	& f' - g'\\
\text{Product Rule}	& fg & f g' + f' g\\
\text{Quotient Rule}	& \frac{f}{g} = fg^{-1} & f (g^{-1})' + f' g^{-1}\\
\text{Chain Rule}	 & f(g(x)) &   f'(g(x))g'(x)\\
\end{array}
$

When a function has a single parameter, $f(x)$, we use $f'$ and $f'(x)$ as shorthands for $\frac{d}{dx} f(x)$.  

You can think of $\frac{d}{dx}$ as an operator that maps a function of one parameter to another function. So to take the derivative with respect to $x$ of $cf(x)$, $\frac{d}{dx} cf(x)$, we are really just moving the constant to the left of the operator: $c\frac{d}{dx}x$ or $c\frac{dx}{dx}$, which makes it more clear that the derivative is $c \times 1 = c$.  We are used to the simplification for the usual arithmetic operators. For example, $-2x = 2 \times -x$ where the negation operator is analogous to our derivative operator. We have just moved to the scalar outside of the operation. Thinking of the derivative as an operator is handy when solving complicated derivatives as we can mechanically apply rules of differentiation.

\section{Review: Vector calculus and partial derivatives}

Real-world problems are rarely so simple that we can express them as a single function of a single parameter, $f(x)$. Let's move on to functions of multiple parameters such as $f(x,y)$. We compute derivatives with respect to one variable (parameter) at a time, giving us two different {\em partial derivatives} for this two-parameter function.  Instead of using operator $\frac{d}{dx}$, we use $\frac{\partial}{\partial x}$ (a stylized $d$ and not the Greek letter $\delta$). For functions of a single parameter, operator $\frac{\partial}{\partial x}$ is equivalent to $\frac{d}{dx}$ but you should  restrict your use of $\frac{d}{dx}$ to functions of a single parameter.

The partial derivative with respect to $x$ is just the usual scalar derivative but treating any other variable in the equation as a constant.  Consider function $f(x,y) = 3x^2y$. The partial derivative with respect to $x$ is written $\frac{\partial}{\partial x} 3x^2y$. There are three constants from the perspective of $\frac{\partial}{\partial x}$: 3, 2, and $y$. Therefore, $\frac{\partial}{\partial x} 3yx^2 = 3y\frac{\partial}{\partial x} x^2 = 3y2x = 6yx$. The partial derivative with respect to $y$ treats $x$ like a constant now: $\frac{\partial}{\partial y} 3x^2y = 3x^2\frac{\partial}{\partial y} y = 3x^2\frac{\partial y}{\partial y} = 3x^2 \times 1 = 3x^2$. 

So, partial derivatives use the usual scalar derivative rules while holding all but a single variable, $x$, as constants for the purpose of computing the partial derivative with respect to $x$.

\section{Matrix calculus}

Moving from vector calculus to matrix calculus now, let's consider what we do with the partial derivatives $\frac{\partial f(x,y)}{\partial x}$ and $\frac{\partial f(x,y)}{\partial y}$ we computed for $f(x,y) = 3x^2y$. Currently they're just kind of floating around and not organized in any way so let's organize them into a horizontal vector:

$[ \frac{\partial f(x,y)}{\partial x}, \frac{\partial f(x,y)}{\partial y}] = [6yx, 3x^2]$

We call this vector the {\em gradient} of $f(x,y)$ and write it as:

$\nabla f(x,y) = [6yx, 3x^2]$.

Now, let's get crazy and consider two functions, both of which take two parameters.  We can keep the same $f(x,y)$ but let's add $g(x,y) = 2x + y^8$.  The gradient for $g$ has two entries, a partial derivative for each parameter. The partials are:

$\frac{\partial g(x,y)}{\partial x} = \frac{\partial 2x}{\partial x} + \frac{\partial y^8}{\partial x} = 2\frac{\partial x}{\partial x} + 0 = 2 \times 1 = 1$

and

$\frac{\partial g(x,y)}{\partial y} = \frac{\partial 2x}{\partial y} + \frac{\partial y^8}{\partial y} = 0 + 8y^7 = 8y^7$

giving us gradient $\nabla g(x,y) = [1, 8y^7]$.

Gradient vectors organize all of the partial derivatives for a specific scalar function. If we have two functions, we can also organize their gradients into a matrix called the {\em Jacobian matrix} where the gradients are rows:

$J =
\begin{bmatrix}
	\nabla f(x,y)\\
	\nabla g(x,y)
\end{bmatrix} = \begin{bmatrix}
	6yx & 3x^2\\
	1 & 8y^7
\end{bmatrix}
$

Welcome to matrix calculus!

\includegraphics[scale=.08]{redbang.png} Before moving on, please note that there are multiple ways to represent the Jacobian. We are using the so-called [numerator layout](https://en.wikipedia.org/wiki/Matrix\_calculus\#Layout\_conventions) but many papers and software will use the transpose of this Jacobian (flip it around its's diagonal).

\subsection{Generalization of the Jacobian}

Assume all vectors are vertical by default of size $n \times 1$. Lowercase letters in bold font such as $\mathbf{x}$ are vectors and math italics font like $x$ or scalars so:

$\mathbf{x} = \begin{bmatrix}
           x_1\\
           x_2\\
           \vdots \\
           x_n\\
           \end{bmatrix}$

Let $\mathbf{y} = \mathbf{f}(\mathbf{x})$ be a vector of $m$ scalar-valued functions, $f_i : R^n \rightarrow R$ with $n=|\mathbf{x}|$. $\mathbf{y} = \mathbf{f}(\mathbf{x})$ is a set of $m$ equations that are functions of $\mathbf{x}$:

$
\begin{array}{lcl}
y_1 & = & f_1(\mathbf{x})\\
y_2 & = & f_2(\mathbf{x})\\
 & \vdots & \\
y_m & = & f_m(\mathbf{x})\\
\end{array}
$

The Jacobian matrix is the collection of all $m \times n$ possible partial derivatives:

$
\frac{\partial \mathbf{y}}{\partial \mathbf{x}} = \begin{bmatrix}
\frac{\partial}{\partial {x}} f_1(\mathbf{x}) \\
\frac{\partial}{\partial {x}} f_2(\mathbf{x})\\
\ldots\\
\frac{\partial}{\partial {x}} f_m(\mathbf{x})
\end{bmatrix}
$

where each $\frac{\partial}{\partial {x}} f_i(\mathbf{x})$ is an $n$-vector because the partial derivative is with respect to a vector, $\mathbf{x}$, whose length is $n = |\mathbf{x}|$.

The width of the Jacobian is the length of the vector $\mathbf{x}$ if we're taking partial derivative with respect to $\mathbf{x}$.  The Jacobian is $m$ rows for $m$ equations.

Show the shape with boxes.  

To represent the trivial equation $\mathbf{y} = \mathbf{x}$, for example, $f_i(\mathbf{x}) = x_i$. In this case, the Jacobian is the identity matrix since $m=n$ (a Jacobian matrix size $m \times n$ where $m = |\mathbf{f}| = n = |\mathbf{x}|$).

$
\frac{\partial \mathbf{y}}{\partial \mathbf{x}} = \begin{bmatrix}
\frac{\partial}{\partial {x}} f_1(\mathbf{x}) \\
\frac{\partial}{\partial {x}} f_2(\mathbf{x})\\
\ldots\\
\frac{\partial}{\partial {x}} f_m(\mathbf{x})
\end{bmatrix} = \begin{bmatrix}
\frac{\partial}{\partial {x_1}} f_1(\mathbf{x}) \frac{\partial}{\partial {x_2}} f_1(\mathbf{x}) \ldots \frac{\partial}{\partial {x_n}} f_1(\mathbf{x}) \\
\frac{\partial}{\partial {x_1}} f_2(\mathbf{x}) \frac{\partial}{\partial {x_2}} f_2(\mathbf{x}) \ldots \frac{\partial}{\partial {x_n}} f_2(\mathbf{x}) \\
\ldots\\
\frac{\partial}{\partial {x_1}} f_m(\mathbf{x}) \frac{\partial}{\partial {x_2}} f_m(\mathbf{x}) \ldots \frac{\partial}{\partial {x_n}} f_m(\mathbf{x}) \\
\end{bmatrix} = \begin{bmatrix}
\frac{\partial}{\partial {x_1}} x_1 \frac{\partial}{\partial {x_2}} x_1 \ldots \frac{\partial}{\partial {x_n}} x_1 \\
\frac{\partial}{\partial {x_1}} x_2 \frac{\partial}{\partial {x_2}} x_2 \ldots \frac{\partial}{\partial {x_n}} x_2 \\
\ldots\\
\frac{\partial}{\partial {x_1}} x_m \frac{\partial}{\partial {x_2}} x_m \ldots \frac{\partial}{\partial {x_n}} x_m \\
\end{bmatrix} = \begin{bmatrix}
           \mathbf{i}_1\\
           \mathbf{i}_2\\
           \vdots \\
           \mathbf{i}_m\\
           \end{bmatrix} = I
$

where $\mathbf{i}_j$ is a vector with $\mathbf{i}_j=1$ and $\mathbf{i}_k=0$ for $k \neq j$ (vector with all zeros except at $j$, which is 1).

It's very important to track whether a matrix is vertical or horizontal, $\mathbf{x}$ or $\mathbf{x}^T$. Also make sure you pay attention to whether something is a scalar-valued function or a vector valued function. Is it $\mathbf{y} = \ldots$ or $y = \ldots$?

The other trick is to reduce everything down to a set of scalar equations and then take all of the partials derivatives. This dictates whether you get a scalar, vector, or matrix Jacobian.

\subsection{Derivatives of vector element-wise binary operators}

Element-wise binary operations on vectors, such as vector addition $\mathbf{w} + \mathbf{x}$, are important because we can express many common vector operations, such as vector-scalar multiplication, as element-wise binary operations.  We can generalize the element-wise binary operations with notation $\mathbf{y} = \mathbf{f(x)} \bigcirc \mathbf{g(x)}$, where we're using $\bigcirc$ to represent any element-wise operator and not the $\circ$ function composition operator.  Here's what that equation looks like when we zoom in to examine the scalar equations:

$\begin{bmatrix}
           y_1\\
           y_2\\
           \vdots \\
           y_n\\
           \end{bmatrix} = \begin{bmatrix}
           f_{1}(\mathbf{x}) \bigcirc g_{1}(\mathbf{x})\\
           f_{n}(\mathbf{x}) \bigcirc g_{2}(\mathbf{x})\\
           \vdots \\
           f_{n}(\mathbf{x}) \bigcirc g_{n}(\mathbf{x})\\
         \end{bmatrix}$

where we write $n$ not $m$ equations vertically to emphasize the fact that the result of element-wise operators give $m=n$ sized vector results.

The general case for the Jacobian is the square matrix:

$
\frac{\partial \mathbf{y}}{\partial \mathbf{x}}  = \begin{bmatrix}
\frac{\partial}{\partial x_1} ( f_{1}(\mathbf{x}) \bigcirc g_{1}(\mathbf{x}) ) & \frac{\partial}{\partial x_2} ( f_{1}(\mathbf{x}) \bigcirc g_{1}(\mathbf{x}) ) & \ldots & \frac{\partial}{\partial x_n} ( f_{1}(\mathbf{x}) \bigcirc g_{1}(\mathbf{x}) )\\\\
\frac{\partial}{\partial x_1} ( f_{2}(\mathbf{x}) \bigcirc g_{2}(\mathbf{x}) ) & \frac{\partial}{\partial x_2} ( f_{2}(\mathbf{x}) \bigcirc g_{2}(\mathbf{x}) ) & \ldots & \frac{\partial}{\partial x_n} ( f_{2}(\mathbf{x}) \bigcirc g_{2}(\mathbf{x}) )\\\\
& \ldots\\\\
\frac{\partial}{\partial x_1} ( f_{n}(\mathbf{x}) \bigcirc g_{n}(\mathbf{x}) ) & \frac{\partial}{\partial x_2} ( f_{n}(\mathbf{x}) \bigcirc g_{n}(\mathbf{x}) ) & \ldots & \frac{\partial}{\partial x_n} ( f_{n}(\mathbf{x}) \bigcirc g_{n}(\mathbf{x}) )\\
\end{bmatrix}
$

That's quite a furball, but the Jacobian is very commonly a diagonal matrix, a matrix that is zero everywhere but the diagonal. 

If $f_i(\mathbf{x})$ and $g_i(\mathbf{x})$ are purely functions of $x_i$, and no other elements of vector $\mathbf{x}$, we can reduce $f_i(\mathbf{x}) \bigcirc g_i(\mathbf{x})$ to $f_i(x_i) \bigcirc g_i(x_i)$. This is technically an abuse of our notation because $f_i$ and $g_i$ are functions of vectors not individual elements. We should really write something like $\hat f_{i}(x_i) = f_{i}(\mathbf{x})$, but that would muddy the equations further (and programmers are comfortable overloading functions) so we'll proceed with the notation anyway.

With the constraint that $f_i(\mathbf{x})$ and $g_i(\mathbf{x})$ can only access $x_i$, the elements along the diagonal of the Jacobian are $\frac{\partial}{\partial x_i} ( f_i(x_i) \bigcirc g_i(x_i) )$:

$
\frac{\partial \mathbf{y}}{\partial \mathbf{x}}  = \begin{bmatrix}
\frac{\partial}{\partial x_1} ( f_{1}(x_1) \bigcirc g_{1}(x_1) )\\
& \frac{\partial}{\partial x_2} (f_{2}(x_2) \bigcirc g_{2}(x_2) ) & &\text{\huge0}\\
& & \ldots \\
\text{\huge0}& & & \frac{\partial}{\partial x_n} (f_{n}(x_n) \bigcirc g_{n}(x_n) )
\end{bmatrix}
$

Or more succinctly,

$\frac{\partial \mathbf{y}}{\partial \mathbf{x}} = diag \left( \frac{\partial}{\partial x_1}(f_{1}(x_1) \bigcirc g_{1}(x_1)),~ \frac{\partial}{\partial x_2}(f_{2}(x_2) \bigcirc g_{2}(x_2)),~ \ldots,~ \frac{\partial}{\partial x_n}(f_{n}(x_n) \bigcirc g_{n}(x_n)) \right)$. 

We know that all off-diagonal entries are 0 because $\frac{\partial}{\partial x_j} ( f_i(x_i) \bigcirc g_i(x_i) ) = 0$ for $j \neq i$. When $j \neq i$, $f_i(x_i) \bigcirc g_i(x_i)$ looks like a constant to the partial derivative operator $\frac{\partial}{\partial x_j}$, and the derivative of any constant is 0.

Because we do lots of simple vector arithmetic, the general function $\mathbf{f(x)}$ is often just the vector $\mathbf{x}$.  Any time the general function is a vector, we know that $f_i(\mathbf{x})$ reduces to $f_i(x_i) = x_i$. For example, vector addition $\mathbf{w + x}$ fits our general case as $\mathbf{f(w)} + \mathbf{g(x)}$ with scalar equations $y_i = f_i(\mathbf{w}) + g_i(\mathbf{x})$. But, that reduces to just $y_i = f_i(w_i) + g_i(x_i) = w_i + x_i$ and so we get partial derivative:

$\frac{\partial}{\partial x_i} ( f_{i}(w_i) + g_{i}(x_i) ) = \frac{\partial}{\partial x_i}(w_i + x_i) = 1 + 0 = 1$

That gives us $\frac{\partial (\mathbf{w+x})}{\partial \mathbf{w}} = I$, the identity matrix, because every element along the diagonal is 1.  Similarly, $\frac{\partial (\mathbf{w+x})}{\partial \mathbf{x}} = I$.

Given the simplicity of this special case, $f_i(\mathbf{x})$ reducing to $f_i(x_i)$, we can easily write out the Jacobian's for the common element-wise binary operations on vectors:

$
\begin{array}{lllll}
\text{Op} & \text{Partial with respect to } \mathbf{w} & \text{Partial with respect to }\mathbf{x}\\
\hline\\

+ & \frac{\partial (\mathbf{w+x})}{\partial \mathbf{w}} = diag(\ldots \frac{\partial (w_i + x_i)}{\partial w_i} \ldots) = diag(\vec{1}) = I & \frac{\partial (\mathbf{w+x})}{\partial \mathbf{x}} =  I\\\\

- & \frac{\partial (\mathbf{w-x})}{\partial \mathbf{w}}  =  diag(\ldots\frac{\partial (w_i - x_i)}{\partial w_i}\ldots) =  diag(\vec{1})  =  I & \frac{\partial (\mathbf{w-x})}{\partial \mathbf{x}}  =  diag(\ldots\frac{\partial (w_i - x_i)}{\partial x_i}\ldots)  =  diag(-\vec{1})  =  -I \\\\

\otimes & \frac{\partial (\mathbf{w \otimes x})}{\partial \mathbf{w}}  =  diag(\ldots\frac{\partial (w_i \times x_i)}{\partial w_i} \ldots)  =  diag(\mathbf{x}) & \frac{\partial (\mathbf{w \otimes x})}{\partial \mathbf{x}}  =  diag(\mathbf{w})\\\\

\oslash & \frac{\partial (\mathbf{w \oslash x})}{\partial \mathbf{w}}  =  diag(\ldots\frac{\partial (w_i / x_i)}{\partial w_i}\ldots)  =  diag(\ldots \frac{1}{x_i} \ldots) & \frac{\partial (\mathbf{w \oslash x})}{\partial \mathbf{x}}  =  diag(\ldots \frac{-w_i}{x_i^2} \ldots)\\

\end{array}
$

\subsection{Derivatives involving scalar expansion}

When we add or multiply scalars to vectors, we are implicitly expanding the scalar to a vector and then performing an element-wise operation. For example, vector-scalar addition 
$\mathbf{y} = \mathbf{x} + c$
is really
$\mathbf{y} = \mathbf{f(x)} + \mathbf{g}(c)$ where $\mathbf{f(x)} = \mathbf{x}$ and $\mathbf{g}(c) = \vec{1} \times c$.
Similarly, vector-scalar multiplication
$\mathbf{y} = \mathbf{x} \times c$
is really
$\mathbf{y} = \mathbf{f(x)} \otimes \mathbf{g}(c)$.

The partial derivative of vector-scalar addition and multiplication with respect to vector $\mathbf{x}$  fall out from our element-wise rule, $\frac{\partial \mathbf{y}}{\partial \mathbf{x}} = diag \left( \ldots \frac{\partial}{\partial x_i} ( f_i(x_i) \bigcirc g_i(c) ) \ldots \right)$, since we've identified $f_i(x_i) = x_i$ and $g_i(c) = c$. Note that functions $\mathbf{f(x)} = \mathbf{x}$ and $\mathbf{g}(c) = \vec{1} \times c$ clearly satisfy our constraint for the special case of diagonal Jacobian (that $f_i(\mathbf{x})$ refer at most to $x_i$ and $g_i(c)$ refer to the $i^{th}$ value of the $\vec{1}$ vector). 

Using the usual rules for scalar partial derivatives, we arrive at the following diagonal elements of the Jacobian for vector-scalar addition:
 
$\frac{\partial}{\partial x_i} ( f_i(x_i) + g_i(c) ) = \frac{\partial (x_i + c)}{\partial x_i} = \frac{\partial x_i}{\partial x_i} + \frac{\partial c}{\partial x_i} = 1 + 0 = 1$

So, $\frac{\partial}{\partial \mathbf{x}} ( \mathbf{x} + c ) = diag(\vec{1}) = I$.

Computing the partial derivative with respect to the single-valued parameter $c$, however, results in a vertical vector not a diagonal matrix whose elements are:
 
$\frac{\partial}{\partial c} ( f_i(x_i) + g_i(c) ) = \frac{\partial (x_i + c)}{\partial c} = \frac{\partial x_i}{\partial c} + \frac{\partial c}{\partial c} = 0 + 1 = 1$

Therefore, $\frac{\partial}{\partial c} ( \mathbf{x} + c ) = \vec{1}$.

The diagonal elements of the Jacobian for vector-scalar multiplication involve the product rule for scalar derivatives:

$\frac{\partial}{\partial x_i} ( f_i(x_i) \otimes g_i(c) ) = x_i \times \frac{\partial c}{\partial x_i} + c \times \frac{\partial x_i}{\partial x_i} = 0 + c = c$

So, $\frac{\partial}{\partial \mathbf{x}} ( \mathbf{x} \times c ) = diag(\vec{1} \times c) = I \times c$. 

The partial derivative with respect to single-valued parameter $c$ is a vertical vector whose elements are:

$\frac{\partial}{\partial c} ( f_i(x_i) \otimes g_i(c) ) = x_i \times \frac{\partial c}{\partial c} + c \times \frac{\partial x_i}{\partial c} = x_i + 0 = x_i$

This gives us $\frac{\partial}{\partial c} ( \mathbf{x} \times c ) = \mathbf{x}$.

\subsection{Derivatives of vector dot product}

For vector dot product $y = \mathbf{f(x)} \cdot \mathbf{g(x)}$ (note $y$ not $\mathbf{y}$ as the result of that product is a scalar), we get:

$
\begin{array}{rcl}
\frac{\partial y}{\partial \mathbf{x}} & = & \mathbf{f(x)} \cdot \frac{\partial \mathbf{g(x)}}{\partial \mathbf{x}} + \mathbf{g(x)} \cdot \frac{\partial \mathbf{f(x)}}{\partial \mathbf{x}}\\\\
 & = & \mathbf{f(x)} \cdot \begin{bmatrix} \frac{\partial g(\mathbf{x})}{\partial x_1},~ \frac{\partial g(\mathbf{x})}{\partial x_2},~ \ldots,~ \frac{\partial g(\mathbf{x})}{\partial x_n} \end{bmatrix}  + \mathbf{g(x)} \cdot \begin{bmatrix} \frac{\partial f(\mathbf{x})}{\partial x_1},~ \frac{\partial f(\mathbf{x})}{\partial x_2},~ \ldots,~ \frac{\partial f(\mathbf{x})}{\partial x_n} \end{bmatrix}\\

\end{array}
$

Note that the partial derivatives are $\frac{\partial f(\mathbf{x})}{\partial x_i}$ not $\frac{\partial f(x_i)}{\partial x_i}$ and not $\frac{\partial f_i(x_i)}{\partial x_i}$.

As an example, consider $y = \mathbf{w} \cdot \mathbf{x}$. The partial derivative with respect to either factor is a gradient, horizontal vector:

$\frac{\partial y}{\partial \mathbf{w}} = \mathbf{w} \cdot \frac{\partial \mathbf{x}}{\partial \mathbf{w}} + \mathbf{x} \cdot \frac{\partial \mathbf{w}}{\partial \mathbf{w}} = \mathbf{w}^T \times \frac{\partial \mathbf{x}}{\partial \mathbf{w}} + \mathbf{x}^T \times \frac{\partial \mathbf{w}}{\partial \mathbf{w}} = \mathbf{w}^T \times \mathbf{0} + \mathbf{x}^T \times I = \mathbf{x}^T$. Similarly, $\frac{\partial y}{\partial \mathbf{x}} = \mathbf{w}^T$.

Note: the dot product transforms into vector-vector multiplication via:

$y = \mathbf{w} \cdot \mathbf{x} = \mathbf{w}^{T} \times \mathbf{x}$

Note: the gradient is a horizontal vector (1 output)!

We can also grind the dot product down into a pure scalar function:

$y = \mathbf{w} \cdot \mathbf{x} = \Sigma_i^n (w_i x_i)$

$\frac{\partial y}{\partial w_j} = \frac{\partial}{\partial w_j} \Sigma_i (w_i x_i) = \Sigma_i \frac{\partial}{\partial w_j} (w_i x_i) = \frac{\partial}{\partial w_j} (w_j x_j) = x_j$

Then:

$\frac{\partial y}{\partial \mathbf{w}} = [ x_1, \ldots, x_n ] = \mathbf{x}^T$

\subsection{Vector sum}

Let $y = sum( \mathbf{f}(\mathbf{x})) = \Sigma_{i=1}^n f_i(\mathbf{x})$ then

$
\begin{array}{lcl}
\frac{\partial y}{\partial \mathbf{x}} & = & \begin{bmatrix} \frac{\partial y}{\partial x_1}, \frac{\partial y}{\partial x_2}, \ldots, \frac{\partial y}{\partial x_n} \end{bmatrix}\\\\
 & = & \begin{bmatrix} \frac{\partial}{\partial x_1} \Sigma_i f_i(\mathbf{x}),~ \frac{\partial}{\partial x_2} \Sigma_i f_i(\mathbf{x}),~ \ldots,~ \frac{\partial}{\partial x_n} \Sigma_i  f_i(\mathbf{x}) \end{bmatrix} \\\\
 & = & \begin{bmatrix} \Sigma_i \frac{\partial f_i(\mathbf{x})}{\partial x_1},~ \Sigma_i \frac{\partial f_i(\mathbf{x})}{\partial x_2},~ \ldots,~ \Sigma_i \frac{\partial f_i(\mathbf{x})}{\partial x_n}  \end{bmatrix}(\text{move derivative inside})\\\\
\end{array}
$

Example: Sum of vector-scalar multiplication.


If $y = sum(\mathbf{x} \times c)$ then $f_i(\mathbf{x},c) = x_i \times c$.

$
\begin{array}{lcl}
\frac{\partial y}{\partial \mathbf{x}} & = & \begin{bmatrix} \Sigma_i \frac{\partial}{\partial x_1} x_i c,~ \Sigma_i \frac{\partial }{\partial x_2} x_i c,~ \ldots,~ \Sigma_i \frac{\partial}{\partial x_n} x_i c  \end{bmatrix}\\\\
 & = & \begin{bmatrix} \frac{\partial}{\partial x_1} x_1 c,~ \frac{\partial }{\partial x_2} x_2 c,~ \ldots,~ \frac{\partial}{\partial x_n} x_n c  \end{bmatrix}\\\\
 & = & \begin{bmatrix} c, c, \ldots, c \end{bmatrix}\\\\
\end{array}
$

The derivative with respect to scalar variable $c$ is $1 \times 1$:

$
\begin{array}{lcl}
\frac{\partial y}{\partial c} & = & \frac{\partial}{\partial c} \Sigma_{i=1}^n (x_i+c)\\
& = & \Sigma_i \frac{\partial}{\partial c} (x_i+c)\\
& = & \Sigma_i (0 + 1)\\
& = & n
\end{array}
$

\subsection{Vector chain rule}

$D[f(g(x))] = f'(g(x))g'(x)$

\section{Acknowledgements}

We thank [Yannet Interian](https://www.usfca.edu/faculty/yannet-interian) (Faculty in MS data science program at University of San Francisco) and [David Uminsky](http://www.cs.usfca.edu/~duminsky/) (Faculty/director of MS data science) for their help with the notation presented here..

\end{document}
